\documentclass{article}
\usepackage[utf8]{inputenc}
\usepackage{amsmath}
\usepackage{amssymb}
\usepackage{tcolorbox}
\usepackage{graphicx}

\begin{document}

\begin{tcolorbox}[colback=blue!10, colframe=blue!50!black, title=AP Calculus Problem]
The table below provides information about the signs of \(f(x)\), \(f'(x)\), and \(f''(x)\) on specific intervals. Use this information to analyze the behavior of \(f(x)\) and answer the questions.

\begin{center}
\begin{tabular}{|c|c|c|c|}
\hline
Interval & \(f'(x)\) & \(f''(x)\) & Behavior of \(f(x)\) \\
\hline
\(x < -1\) & \(+\) (positive) & \(+\) (positive) & Increasing and concave up \\
\(-1 < x < 1\) & \(-\) (negative) & \(+\) (positive) & Decreasing and concave up \\
\(x = 1\) & \(-\) (negative) & \(0\) & No extrema, inflection point \\
\(1 < x < 3\) & \(-\) (negative) & \(-\) (negative) & Decreasing and concave down \\
\(x = 3\) &  (undefined) & \(0\) & No extrema, inflection point \\
\(x > 3\) & \(-\) (negative) & \(-\) (negative) & Increasing and concave down \\
\hline
\end{tabular}
\end{center}

\textbf{Important Notes:}
\begin{itemize}
    \item \(f'(x) \neq 0\) on the entire domain. Therefore, there are no critical points.
    \item Inflection points occur at \(x = 1\) and \(x = 3\) because \(f''(x) = 0\) and the concavity changes at these points.
\end{itemize}

\textbf{Questions:}
\begin{enumerate}
    \item Identify all intervals where \(f(x)\) is increasing or decreasing. Justify your answer.
    \item Identify all intervals where \(f(x)\) is concave up or concave down. Justify your answer.
    \item Explain why there are no local maxima or minima for \(f(x)\).
    \item Sketch the graph of \(f(x)\) based on the given information.
    \item Does \(f(x)\) have an absolute maximum, absolute minimum, or neither? Justify your answer.
\end{enumerate}
\end{tcolorbox}

\textbf{Answers:}
\begin{enumerate}
    \item \textbf{Intervals of Increase/Decrease:}
    \begin{itemize}
        \item \(f(x)\) is increasing on \(x < -1\) and \(x > 3\), because \(f'(x) > 0\) on these intervals.
        \item \(f(x)\) is decreasing on \(-1 < x < 3\), because \(f'(x) < 0\) on this interval.
    \end{itemize}
    \item \textbf{Concavity:}
    \begin{itemize}
        \item \(f(x)\) is concave up on \(x < 1\), because \(f''(x) > 0\) on this interval.
        \item \(f(x)\) is concave down on \(x > 1\), because \(f''(x) < 0\) on this interval.
    \end{itemize}
    \item \textbf{Explanation of No Local Extrema:}
    \begin{itemize}
        \item \(f'(x) \neq 0\) on the entire domain, which means there are no critical points. Since a critical point is necessary for local maxima or minima, \(f(x)\) does not have any local extrema.
        \item Inflection points at \(x = 1\) and \(x = 3\) do not correspond to extrema, as they only indicate a change in concavity.
    \end{itemize}
    \item \textbf{Graph Sketch:}
    \begin{itemize}
        \item The graph of \(f(x)\) should show increasing behavior on \(x < -1\) and \(x > 3\), with concave up on \(x < 1\) and concave down on \(x > 1\).
        \item The graph should decrease on \(-1 < x < 3\) with a change in concavity at \(x = 1\) and \(x = 3\).
    \end{itemize}
    \item \textbf{Absolute Extrema:}
    \begin{itemize}
        \item \(f(x)\) does not have an absolute maximum or minimum because the function extends indefinitely on \((- \infty, \infty)\).
    \end{itemize}
\end{enumerate}



\newpage
\begin{tcolorbox}[colback=blue!10, colframe=blue!50!black, title=AP Calculus Problem]
Let \(y = f(x)\) be the solution to the differential equation
\[
\frac{dy}{dx} = 2x + 3,
\]
with the initial condition \(f(1) = 5\).

\textbf{Part A:} Find the equation of the tangent line to the solution curve at the point where \(x = 1\).

\textbf{Part B:} Solve the differential equation to find the particular solution \(y = f(x)\) that satisfies the given initial condition.
\end{tcolorbox}

\textbf{Solutions:}

\textbf{Part A: Finding the equation of the tangent line}

The equation of a tangent line is given by:
\[
y - y_1 = m(x - x_1),
\]
where \(m\) is the slope of the tangent line and \((x_1, y_1)\) is a point on the curve.

The slope of the tangent line is the value of \(\frac{dy}{dx}\) at \(x = 1\):
\[
\frac{dy}{dx} = 2x + 3.
\]
Substitute \(x = 1\):
\[
m = 2(1) + 3 = 5.
\]

At \(x = 1\), \(y = f(1) = 5\). Therefore, the point on the curve is \((1, 5)\).

The equation of the tangent line is:
\[
y - 5 = 5(x - 1),
\]
or equivalently:
\[
y = 5x.
\]

\textbf{Part B: Solving the differential equation}

The differential equation is:
\[
\frac{dy}{dx} = 2x + 3.
\]

To solve, integrate both sides with respect to \(x\):
\[
y = \int (2x + 3) \, dx.
\]

Using the power rule:
\[
y = x^2 + 3x + C,
\]
where \(C\) is the constant of integration.

Using the initial condition \(f(1) = 5\), substitute \(x = 1\) and \(y = 5\) into the equation:
\[
5 = (1)^2 + 3(1) + C.
\]
Simplify:
\[
5 = 1 + 3 + C \implies C = 1.
\]

Thus, the particular solution is:
\[
y = x^2 + 3x + 1.
\]

\textbf{Final Answer:}
\begin{itemize}
    \item \textbf{Part A:} The equation of the tangent line is \(y = 5x\).
    \item \textbf{Part B:} The particular solution to the differential equation is \(y = x^2 + 3x + 1\).
\end{itemize}

\newpage
\begin{tcolorbox}[colback=blue!10, colframe=blue!50!black, title=AP Calculus Problem]
Let \(y = f(x)\) be the solution to the differential equation
\[
\frac{dy}{dx} = y \ln(x),
\]
with the initial condition \(f(e) = 4\).

\textbf{Part A:} Find the equation of the tangent line to the solution curve at the point where \(x = e\).

\textbf{Part B:} Solve the differential equation to find the particular solution \(y = f(x)\) that satisfies the given initial condition.

\textbf{Part C:} Determine the behavior of the solution as \(x \to \infty\) and as \(x \to 0^+\).
\end{tcolorbox}

\textbf{Solutions:}

\textbf{Part A: Finding the equation of the tangent line}

The equation of a tangent line is given by:
\[
y - y_1 = m(x - x_1),
\]
where \(m\) is the slope of the tangent line and \((x_1, y_1)\) is a point on the curve.

The slope of the tangent line is the value of \(\frac{dy}{dx}\) at \(x = e\). Using the differential equation:
\[
\frac{dy}{dx} = y \ln(x).
\]
Substitute \(x = e\) and \(y = f(e) = 4\):
\[
\frac{dy}{dx} = 4 \ln(e) = 4(1) = 4.
\]

The point on the curve is \((e, 4)\). The equation of the tangent line is:
\[
y - 4 = 4(x - e),
\]
or equivalently:
\[
y = 4x - 4e + 4.
\]

---

\textbf{Part B: Solving the differential equation}

The differential equation is:
\[
\frac{dy}{dx} = y \ln(x).
\]

This is a separable differential equation. Rewrite it as:
\[
\frac{1}{y} \, dy = \ln(x) \, dx.
\]

Integrate both sides:
\[
\int \frac{1}{y} \, dy = \int \ln(x) \, dx.
\]

The left-hand side integrates to:
\[
\ln|y| + C_1.
\]

The right-hand side requires integration by parts. Let \(u = \ln(x)\) and \(dv = dx\), so that \(du = \frac{1}{x} dx\) and \(v = x\). Then:
\[
\int \ln(x) \, dx = x \ln(x) - \int x \cdot \frac{1}{x} \, dx = x \ln(x) - x + C_2.
\]

Thus, the equation becomes:
\[
\ln|y| = x \ln(x) - x + C.
\]

Exponentiate both sides to solve for \(y\):
\[
y = e^{x \ln(x) - x + C}.
\]

Simplify the exponent:
\[
y = e^C \cdot e^{x \ln(x) - x}.
\]

Let \(A = e^C\) (a constant), so:
\[
y = A e^{x \ln(x) - x}.
\]

Use the initial condition \(f(e) = 4\) to find \(A\). Substitute \(x = e\) and \(y = 4\):
\[
4 = A e^{e \ln(e) - e}.
\]
Simplify:
\[
4 = A e^{e(1) - e} = A e^0 = A.
\]

Thus, \(A = 4\), and the particular solution is:
\[
y = 4 e^{x \ln(x) - x}.
\]

---

\textbf{Part C: Behavior of the solution}

1. As \(x \to \infty\):
   \[
   y = 4 e^{x \ln(x) - x}.
   \]
   The exponent \(x \ln(x) - x = x(\ln(x) - 1)\). Since \(\ln(x)\) grows slower than \(x\), \(\ln(x) - 1 > 0\) for large \(x\). Therefore, \(y \to \infty\) as \(x \to \infty\).

2. As \(x \to 0^+\):
   \[
   y = 4 e^{x \ln(x) - x}.
   \]
   For \(x \to 0^+\), \(\ln(x) \to -\infty\), so the term \(x \ln(x)\) dominates. Since \(x \ln(x) \to 0\) as \(x \to 0^+\), the solution approaches \(y \to 0^+\).

---

\textbf{Final Answers:}
\begin{itemize}
    \item \textbf{Part A:} The equation of the tangent line is \(y = 4x - 4e + 4\).
    \item \textbf{Part B:} The particular solution to the differential equation is \(y = 4 e^{x \ln(x) - x}\).
    \item \textbf{Part C:} 
    \begin{itemize}
        \item As \(x \to \infty\), \(y \to \infty\).
        \item As \(x \to 0^+\), \(y \to 0^+\).
    \end{itemize}
\end{itemize}


\newpage

\begin{tcolorbox}[colback=blue!10, colframe=blue!50!black, title=Riemann Sum Problem]
A function \(f(x)\) has the values given in the table below for \(x\) and \(f(x)\). Use the subintervals from the table to approximate the area bounded by the \(x\)-axis and the curve, on \([0, 12]\). Be sure to include a diagram for each method.

\begin{center}
\begin{tabular}{|c|c|c|c|c|c|c|c|}
\hline
\(x\) & 0 & 2 & 4 & 6 & 8 & 10 & 12 \\
\hline
\(f(x)\) & 5 & 7 & 6 & 4 & 2 & 1 & 3 \\
\hline
\end{tabular}
\end{center}

\textbf{Part A:} Approximate the area using the \textbf{Left Hand Riemann Sum}. 

\textbf{Part B:} Approximate the area using the \textbf{Right Hand Riemann Sum}.

\textbf{Part C:} Based on your results in Part A and Part B, describe what the Riemann sum is approximating in the context of this problem.
\end{tcolorbox}

\textbf{Solutions:}

\textbf{Part A: Left Hand Riemann Sum}

The Left Hand Riemann Sum uses the left endpoint of each subinterval. The subinterval widths are all \(\Delta x = 2\), and the left-hand values of \(f(x)\) are \(f(0)\), \(f(2)\), \(f(4)\), \(f(6)\), \(f(8)\), and \(f(10)\).

\[
A = \Delta x \cdot \big(f(0) + f(2) + f(4) + f(6) + f(8)\big)
\]

Substitute the values:
\[
A = 2 \cdot \big(5 + 7 + 6 + 4 + 2\big) = 2 \cdot 24 = 48.
\]

Thus, the Left Hand Riemann Sum is:
\[
A \approx 48.
\]

---

\textbf{Part B: Right Hand Riemann Sum}

The Right Hand Riemann Sum uses the right endpoint of each subinterval. The subinterval widths are again \(\Delta x = 2\), and the right-hand values of \(f(x)\) are \(f(2)\), \(f(4)\), \(f(6)\), \(f(8)\), \(f(10)\), and \(f(12)\).

\[
A = \Delta x \cdot \big(f(2) + f(4) + f(6) + f(8) + f(10) + f(12)\big)
\]

Substitute the values:
\[
A = 2 \cdot \big(7 + 6 + 4 + 2 + 1 + 3\big) = 2 \cdot 23 = 46.
\]

Thus, the Right Hand Riemann Sum is:
\[
A \approx 46.
\]

---

\textbf{Part C: Interpretation}

The Riemann sums approximate the area under the curve \(f(x)\) from \(x = 0\) to \(x = 12\). This represents the total "accumulated quantity" described by \(f(x)\) over the interval \([0, 12]\). For example, if \(f(x)\) represents a rate of change (e.g., velocity), the Riemann sum would approximate the total change (e.g., total distance) over the interval.

The difference between the Left Hand and Right Hand Riemann Sums arises because the function \(f(x)\) is not constant over the subintervals. If we increase the number of subintervals (decrease \(\Delta x\)), the approximation will become more accurate.




\end{document}
